\chapter{Introduction}
\label{ch:introduction}

\section{Background}

Robotics is seen as an important technology for the future of agriculture, for example to address the shortage of workers and to reduce environmental impacts such as soil compaction and the use of pesticides~\citep{duckett2018}. Typical applications of agricultural robots are weeding, harvesting and the monitoring of crops~\citep{oliveira2021}. Autonomous mobile robots can take over such repetitive tasks. Before a robot can do any of these tasks, it has to know its surroundings and be able to move in them safely without hitting obstacles. This project therefore focuses on these basic capabilities, mapping and navigation, on which later tasks can build.

This project builds on a previous project at the EDXI Lab of ZHAW, which looked at an autonomous robot for plant management on green rooftops with photovoltaic installations~\citep{graham2024}. In that project, a robot chassis, a basic control logic and design specifications were developed, and a digital twin of the robot was built with ROS2. The real robot did not drive autonomously yet: it had no sensors for the environment and no wheel encoders, so odometry, mapping and navigation were not possible. The previous work also recommended using a long-term support release of ROS2 instead of the unstable Rolling distribution.

For this project, the EDXI Lab provided a commercial robot platform, a Waveshare RaspRover, which already includes wheel encoders, an IMU and a motor controller, together with an RPLIDAR A2M12 lidar sensor.

\section{Objectives}

The objective of this project is to implement robot control for movement and environmental mapping, so that the robot can map and navigate the Schlosshof of ZHAW Wädenswil using ROS2, with all processing done on the robot itself (edge computing). Two results are expected:

\begin{enumerate}
  \item \textbf{Robot and sensor control}
  \item \textbf{Collision avoidance and mapping} with the lidar sensor
\end{enumerate}

The work follows the guiding questions from the project outline (Tab.~\ref{tab:questions}).

\begin{table}[H]
  \centering
  \small
  \caption{Guiding questions of the project and the chapters in which they are answered.}
  \label{tab:questions}
  \begin{tabularx}{\textwidth}{@{}l X l@{}}
    \toprule
    \textbf{Phase} & \textbf{Guiding question} & \textbf{Chapter} \\
    \midrule
    Preparation        & What approaches exist for mapping and navigation of mobile robots? & \ref{ch:theory} \\
                       & What are the challenges of self-driving robots? & \ref{ch:theory} \\
    \midrule
    Setup of robot     & How does the robot move? & \ref{ch:methods}, \ref{ch:results} \\
    and sensors        & Which sensors and actuators does the robot have? & \ref{ch:methods} \\
                       & How do the robot's components communicate with each other? & \ref{ch:methods} \\
    \midrule
    Mapping and        & How can the robot capture and map its environment? & \ref{ch:methods}, \ref{ch:results} \\
    navigation         & How does the robot determine its position on the map? & \ref{ch:methods}, \ref{ch:results} \\
                       & How does the robot detect and avoid obstacles? & \ref{ch:methods}, \ref{ch:results} \\
                       & How can the robot reach a given position? & \ref{ch:methods}, \ref{ch:results} \\
    \midrule
    Evaluation         & How accurate is the generated map? & \ref{ch:results}, \ref{ch:discussion} \\
                       & How reliably does the robot reach its target? & \ref{ch:results}, \ref{ch:discussion} \\
                       & What are the limitations of the system? & \ref{ch:discussion} \\
    \bottomrule
  \end{tabularx}
\end{table}

\section{Scope}

The project only covers the perception of the environment and the movement of the robot. Plant detection, additional actuators and the robotic arm that is developed for the robot by students in Madrid are not part of this work.